\documentclass[11pt,a4paper]{article}
\usepackage[margin=2.5cm]{geometry}
\usepackage{amsmath,amssymb}
\usepackage{booktabs}
\usepackage{hyperref}
\usepackage{listings}
\lstset{basicstyle=\ttfamily\small,columns=fullflexible,frame=single,language=Fortran}

\newcommand{\bv}{\mathbf{v}}
\newcommand{\bx}{\mathbf{x}}
\newcommand{\kms}{\mathrm{km\,s^{-1}}}
\newcommand{\Msun}{\mathrm{M_\odot}}
\newcommand{\vs}{v_{\mathrm s}}
\newcommand{\avg}[1]{\left\langle #1 \right\rangle}

\title{Velocities that follow Larson's size--linewidth relation\\
in the HENDOS fractal generator}
\author{Sheffield star formation and dynamics group}
\date{\today}

\begin{document}
\maketitle

\section{Overview}

The HENDOS initial conditions generator \texttt{makeics} can make a box
fractal of stars (Goodwin \& Whitworth 2004) and give the stars
velocities in one of four ways. One of them, \texttt{larson}
(made by the routine \texttt{larsonfractal}), is designed so that the
velocity difference between two stars grows as the square root of their
separation,
\begin{equation}
  \delta v \propto d^{1/2},
  \label{eq:aim}
\end{equation}
which is the slope of the size--linewidth relation of molecular clouds.
This note explains how that is done, shows that the result follows
equation~(\ref{eq:aim}) exactly on average, describes how the velocities
are put on the observed relation in $\kms$, and compares the other three
velocity options.

The short answer: \texttt{larson} gives the right \emph{slope} at every
scale by construction, and the \emph{normalisation} is set to the
observed relation of Solomon et al.\ (1987). Unlike every other option
in the generator, the velocities are therefore not scaled to a virial
ratio: the virial ratio is an outcome. The other options do not give
this slope.

\section{The size--linewidth relation}

Larson (1981) found that the velocity dispersion $\sigma$ of molecular
clouds and of regions within them increases with their size $L$ as
\begin{equation}
  \sigma \simeq 1.10\,\kms\left(\frac{L}{\mathrm{pc}}\right)^{0.38}.
\end{equation}
Later work with more uniform data found a steeper slope; the relation
most often quoted is that of Solomon et al.\ (1987),
\begin{equation}
  \sigma \simeq 0.72\,\kms\left(\frac{R}{\mathrm{pc}}\right)^{0.5},
  \label{eq:srby}
\end{equation}
where $\sigma$ is the one-dimensional velocity dispersion and $R$ the
radius of the cloud. Heyer et al.\ (2009) showed that the coefficient is
not universal but increases with the surface density of the cloud. The
generator aims at the slope $1/2$ of equation~(\ref{eq:srby}).

The relation is a statement about gas. Using it for stars assumes that
stars form with the velocities of the gas they form from, so that stars
that form close together have similar velocities, and stars that form
far apart differ by more.

\section{The method}

\subsection{A Gaussian velocity field with a linear covariance}

Once the $N$ star positions $\bx_i$ are fixed, each Cartesian component
of the velocities is drawn as a set of $N$ Gaussian random numbers with
zero mean and covariance
\begin{equation}
  \avg{v_{i,a}\,v_{j,b}} = \delta_{ab}\,C_{ij},
  \qquad
  C_{ij} = \vs^2\left(D - \tfrac12 d_{ij}\right),
  \label{eq:cov}
\end{equation}
where $a,b$ label the components, $d_{ij} = |\bx_i-\bx_j|$ is the
separation of stars $i$ and $j$, $\vs$ is a velocity scale
(\texttt{vscale}, with units of velocity per square root of length) and
$D$ is a constant length. The code uses $D = 10\,d_{\max}$, ten times
the largest separation in the system. The three components are
independent of each other.

\subsection{Why this gives the square-root law}

The mean square difference between one component of the velocities of
two stars is
\begin{equation}
  \avg{(v_{i,a}-v_{j,a})^2} = C_{ii} + C_{jj} - 2C_{ij}
  = \vs^2 D + \vs^2 D - 2\vs^2\left(D - \tfrac12 d_{ij}\right)
  = \vs^2\, d_{ij} .
\end{equation}
The constant $D$ cancels, and what is left is exactly proportional to
the separation. Summing the three components,
\begin{equation}
  \avg{|\bv_i-\bv_j|^2} = 3\,\vs^2\,d_{ij},
  \qquad\text{so}\qquad
  \delta v_{\mathrm{rms}}(d) = \sqrt{3}\,\vs\,d^{1/2}.
  \label{eq:sf}
\end{equation}
This holds for every pair of stars, at every separation, whatever the
positions are. It does not depend on the fractal dimension.

\subsection{The velocity dispersion of a region}

Larson's relation is usually stated for the velocity dispersion of a
region rather than for pairs. The two are connected by an identity: for
any set $S$ of $n$ stars with mean velocity $\bar\bv_S$,
\begin{equation}
  \sigma_S^2 \equiv \frac1n\sum_{i\in S}|\bv_i-\bar\bv_S|^2
  = \frac{1}{2n^2}\sum_{i,j\in S}|\bv_i-\bv_j|^2 .
\end{equation}
Taking the average with equation~(\ref{eq:sf}),
\begin{equation}
  \avg{\sigma_S^2} = \tfrac32\,\vs^2\,\bar d_S,
  \qquad
  \avg{\sigma_{S,\mathrm{1D}}^2} = \tfrac12\,\vs^2\,\bar d_S,
  \label{eq:region}
\end{equation}
where $\bar d_S$ is the mean separation of pairs of stars in the region.
The mean separation is proportional to the size of the region, so
$\sigma \propto L^{1/2}$. For example, the mean distance between two
points in a uniform sphere of radius $R$ is $36R/35$, so for stars spread
uniformly through such a sphere
\begin{equation}
  \sigma_{\mathrm{1D}} = \left(\frac{18}{35}\right)^{1/2}\vs\,R^{1/2}
  \simeq 0.72\,\vs\,R^{1/2}.
\end{equation}
By coincidence the numerical factor is the same as the coefficient in
equation~(\ref{eq:srby}): a velocity scale $\vs =
1\,\kms\,\mathrm{pc^{-1/2}}$ corresponds to the Solomon et al.\ relation
for a uniformly filled region.

\subsection{The constant length and the bulk velocity}

Equation~(\ref{eq:cov}) gives every star a velocity variance $\vs^2 D$
per component, which is much larger than the differences between
stars. Almost all of it is a single velocity shared by every star: the
correlation between any two stars is at least $1 - d_{\max}/2D = 0.95$.
This shared velocity carries no information and is removed when the
generator moves the stars to their centre-of-velocity frame. What
remains are the differences, which obey equation~(\ref{eq:sf}).

$D$ is needed to make equation~(\ref{eq:cov}) a valid covariance, that
is, to make the matrix $C$ positive definite so that Gaussian numbers
with that covariance exist. For weights $w_i$,
\begin{equation}
  \sum_{i,j} w_i C_{ij} w_j = \vs^2 D\Big(\sum_i w_i\Big)^2
  - \tfrac12\vs^2\sum_{i,j} w_i\,d_{ij}\,w_j .
\end{equation}
For distinct points in Euclidean space the second sum is negative
whenever $\sum_i w_i = 0$ (the distance $|\bx-\mathbf{y}|$ is
conditionally negative definite; Schoenberg 1938), so the form is
positive in those directions, and a large enough $D$ makes it positive
in the remaining direction. In the tests of Section~\ref{sec:tests} the
smallest eigenvalue of $C$ was positive in every case with $D =
10\,d_{\max}$.

\subsection{Relation to turbulent velocity fields}

A Gaussian field whose increments have variance proportional to
$d^{2H}$ is a fractional Brownian field with Hurst exponent $H$.
Equation~(\ref{eq:sf}) is the case $H = 1/2$, for which the
three-dimensional power spectrum of the velocity is $P(k) \propto
k^{-(3+2H)} = k^{-4}$. This is the spectrum commonly used to set up
turbulent velocity fields in simulations of star formation, and is that
of Burgers (shock-dominated) turbulence. The method is therefore
equivalent to sampling such a turbulent field at the positions of the
stars, without having to make the field on a grid. Because the three
components are independent, the field is a mixture of solenoidal and
compressive motions in the natural ratio $2:1$.

\section{Implementation}

The routine is \texttt{larsonfractal} in
\texttt{HENDOS/src/ics/fractalsubs.f90}. It
\begin{enumerate}
\item calls the routine \texttt{makefractal} to make the positions and
  cut the number of stars to $N$;
\item finds the largest separation $d_{\max}$ and sets $D =
  10\,d_{\max}$;
\item fills the $N\times N$ matrix $C$ of equation~(\ref{eq:cov});
\item factorises it as $C = \mathsf{L}\,\mathsf{\Delta}\,
  \mathsf{L}^{\mathsf T}$, with $\mathsf L$ unit lower triangular and
  $\mathsf\Delta$ diagonal (a Cholesky decomposition without square
  roots);
\item for each component draws $N$ independent Gaussian numbers $z_k$ of
  unit variance and sets
  \begin{equation}
    v_i = \sum_k \mathsf{L}_{ik}\,\Delta_k^{1/2}\,z_k ,
  \end{equation}
  which has covariance $\mathsf{L}\mathsf{\Delta}\mathsf{L}^{\mathsf T} =
  C$ as required;
\item subtracts the mean velocity.
\end{enumerate}
The Gaussian numbers come from \texttt{gasdev}, which uses the random
number generator \texttt{ran3}.

\paragraph{Cost.} The matrix needs $8N^2$ bytes (8\,MB for $N=10^3$,
0.8\,GB for $N=10^4$) and the factorisation takes of order $N^3/3$
operations. The method is practical for a few thousand stars and
becomes slow beyond about $10^4$.

\section{Normalisation}
\label{sec:norm}

For every other kind of initial conditions \texttt{makeics} multiplies
all the velocities by one factor to give the virial ratio the user asks
for. A Larson fractal is the exception: once the stars have their masses
and the fractal has its size, the velocities are set to the observed
relation in $\kms$, and the virial ratio \texttt{qvir} in the input is
not used.

\subsection{Setting the coefficient}

The routine \texttt{larsonfractal} is called with radius 1 and $\vs =
1$, so that by equation~(\ref{eq:sf}) the mean square difference of one
velocity component between two stars equals their separation in units of
the radius. The positions are then multiplied by the radius $R$ (in pc)
and the velocities by
\begin{equation}
  \vs\,R^{1/2}, \qquad \vs = \texttt{vlarson},
\end{equation}
which gives, with $d$ in pc,
\begin{equation}
  \avg{(v_{i,a}-v_{j,a})^2} = \vs^2\,d
  \quad(\kms)^2 .
\end{equation}
The input \texttt{vlarson} is therefore the rms difference, in $\kms$, of
one velocity component between two stars 1\,pc apart. Its standard value
is $1\,\kms$, for which equation~(\ref{eq:region}) gives a
one-dimensional velocity dispersion of $0.72\,(R/\mathrm{pc})^{1/2}\,
\kms$ for a region filled uniformly to radius $R$: the relation of
Solomon et al.\ (1987), equation~(\ref{eq:srby}). Other coefficients
(for example those that depend on surface density; Heyer et al.\ 2009)
are obtained by changing \texttt{vlarson}. Larson's (1981) own relation
has a different slope and cannot be matched in this way.

A fractal is not uniformly filled, so the dispersion of the whole
system is not exactly $0.72\,\vs R^{1/2}$: it is $(\bar d/2)^{1/2}\vs$,
where $\bar d$ is the mean separation of the stars, which is smaller
for a clumpier fractal. What is fixed is the relation between pairs of
stars.

\subsection{The virial ratio that results}

From equation~(\ref{eq:region}) applied to the whole system of mass
$M$, the kinetic energy is $T = \tfrac34 M\vs^2\bar d$ on average, so the
virial ratio is
\begin{equation}
  Q = \frac{T}{|\Omega|} = \frac{3\,M\,\vs^2\,\bar d}{4\,|\Omega|} .
  \label{eq:q}
\end{equation}
For a uniform sphere ($|\Omega| = 3GM^2/5R$, $\bar d = 36R/35$) this is
\begin{equation}
  Q = \frac97\,\frac{\vs^2 R^2}{GM}
  \simeq 0.30\left(\frac{\vs}{1\,\kms}\right)^2
  \left(\frac{R}{\mathrm{pc}}\right)^2
  \left(\frac{M}{1000\,\Msun}\right)^{-1} .
\end{equation}
The virial ratio rises as $R^2/M$: large, low-mass systems are unbound
and small, massive ones are cold. A clumpy fractal has a larger
$|\Omega|$ and a smaller $\bar d$ than a uniform sphere of the same mass
and radius, so a smaller $Q$. Table~\ref{tab:q} gives the values found
for the standard mass function. \texttt{makeics} prints the virial ratio
of the system it has made and records it in the file.

\begin{table}[h]
\centering
\begin{tabular}{cccc}
\toprule
$R$ (pc) & $M$ ($\Msun$) & $Q$ & $\sigma_{\mathrm{1D}}$ ($\kms$) \\
\midrule
0.5 & 377 & $0.10 \pm 0.04$ & 0.46 \\
1.0 & 368 & $0.43 \pm 0.16$ & 0.68 \\
2.0 & 375 & $1.54 \pm 0.80$ & 0.95 \\
\bottomrule
\end{tabular}
\caption{Virial ratio $Q$ and one-dimensional velocity dispersion of
Larson fractals of $N=1000$ stars, fractal dimension 1.6, with
$\texttt{vlarson} = 1\,\kms$ and masses from the standard mass function.
Means over six fractals; the error is the scatter between fractals.}
\label{tab:q}
\end{table}

The scatter in $Q$ between fractals is large, because both $|\Omega|$
and the few large-scale motions that carry most of the kinetic energy
vary from one fractal to the next. If a particular virial ratio is
needed as well as the Larson slope, the radius or \texttt{vlarson} can
be adjusted until the printed value is the one wanted.

\section{The other velocity options}

The generator has three other velocity fields, selected by
\texttt{velocities} in the input file and made by the routine
\texttt{makefractal} (the \texttt{vtype} given below is its argument).
For these the velocities are scaled to the virial ratio asked for. They are built into the fractal as
it is made: at each level of the hierarchy a box (parent) is divided
into eight (with the standard two boxes per side) and some of the
children survive. Let $\delta_\ell$ be the half-width of a box at level
$\ell$, which halves at each level.

\paragraph{\texttt{inherit} (\texttt{vtype='f'}).} Each child takes its
parent's velocity plus a Gaussian random velocity of dispersion
$\vs\,\delta_\ell$ per component. Two stars whose last common ancestor
is at level $\ell$ share all the kicks down to that level and differ by
those below it, so
\begin{equation}
  \avg{|\bv_i-\bv_j|^2} = 6\,\vs^2\sum_{k>\ell}\delta_k^2
  = 6\,\vs^2\delta_\ell^2\sum_{m\ge1}4^{-m} = 2\,\vs^2\delta_\ell^2 .
\end{equation}
Stars that separate at level $\ell$ are typically a box width apart, so
this is $\delta v \propto d$, a slope of 1 rather than $1/2$. The
measured slope is somewhat less than 1 (Table~\ref{tab:slopes}) because
separation is only loosely tied to the level: stars in neighbouring
boxes can be close together in space but have a distant common
ancestor. A slope of $1/2$ could be made this way by using kicks
proportional to $\delta_\ell^{1/2}$; this is not implemented.

\paragraph{\texttt{standard} (\texttt{vtype='s'}).} The velocities of
the original \texttt{makefractal}: each child takes its parent's
velocity plus a Gaussian random velocity of the \emph{same} dispersion
at every level. Then $\avg{|\bv_i-\bv_j|^2} \propto \ell_{\max}-\ell$,
the number of levels below the common ancestor, which grows only as the
logarithm of the separation. Nearby stars do have similar velocities,
but there is no power law, and a fitted slope is small.

\paragraph{\texttt{random} (\texttt{vtype='u'}).} Every star has an
independent Gaussian velocity. The velocity difference does not depend
on separation.

\section{Tests}
\label{sec:tests}

The tests were made with a line-by-line Python copy of the routine, not
with the Fortran, which had not been compiled when this was written.
Fractals of $N = 1000$ stars with fractal dimension 1.6 and radius 1
were made, the rms velocity difference of all pairs of stars was
measured in 12 logarithmic bins of separation (from about 0.05 to 1.5),
and a power law was fitted.

\begin{table}[h]
\centering
\begin{tabular}{llcc}
\toprule
\texttt{velocities} & \texttt{vtype} & slope, \texttt{coresize} $=0$ &
  slope, \texttt{coresize} $=0.003$ \\
\midrule
\texttt{larson}   & -- & $0.50 \pm 0.07$ & $0.50 \pm 0.08$ \\
\texttt{inherit}  & \texttt{f} & $0.80 \pm 0.09$ & $0.79 \pm 0.12$ \\
\texttt{standard} & \texttt{s} & $0.12 \pm 0.02$ & $0.10 \pm 0.02$ \\
\texttt{random}   & \texttt{u} & $0.00 \pm 0.00$ & $0.00 \pm 0.00$ \\
\bottomrule
\end{tabular}
\caption{Slope of the rms velocity difference between stars against
their separation. Each entry is the mean over ten fractals; the error is
the scatter between fractals, not the error on the mean.}
\label{tab:slopes}
\end{table}

Four further checks were made of \texttt{larson}:
\begin{itemize}
\item Drawing 200 sets of velocities for one set of positions, the
  ratio $\avg{|\bv_i-\bv_j|^2}/3\vs^2 d$ was between 0.99 and 1.04 in
  every bin, confirming equation~(\ref{eq:sf}).
\item The matrix $C$ was positive definite for every fractal made.
\item For fractal dimensions 2.0 and 2.6 the slopes of single fractals
  were $0.46\pm0.05$ and $0.44\pm0.05$ (five fractals each).
\item With the velocities set as in Section~\ref{sec:norm} and
  $\texttt{vlarson} = 1\,\kms$, the coefficient measured from the pairs
  of stars was 1.02 to 1.04, and uniform spheres (fractal dimension 3)
  of radius 1\,pc had $\sigma_{\mathrm{1D}} = 0.71\,\kms$, against the
  0.72 expected.
\end{itemize}

The scatter between fractals is real and is worth remembering: any one
set of initial conditions is a single draw, and its largest scales are
sampled by only a few independent motions, so its measured slope can
differ from $1/2$ by about 0.1. Fitting the logarithm of a noisy
average also biases the slope of a single fractal slightly low. The
relation is exact only on average.

\section{Limitations}

\begin{itemize}
\item The virial ratio is not controlled: it follows from the mass and
  size (Section~\ref{sec:norm}) and varies a good deal between fractals.
\item The observed coefficient is not universal (Heyer et al.\ 2009); the
  standard value is that of Solomon et al.\ (1987).
\item The velocities are Gaussian. Real turbulence is intermittent, with
  non-Gaussian tails, and has correlations between velocity and density
  that are absent here: the velocity field knows about the positions of
  the stars only through their separations.
\item The observed relation is for gas, measured from line widths along
  the line of sight within projected regions. The generator imposes the
  equivalent relation between stars in three dimensions.
\item Memory and time limit the method to about $10^4$ stars.
\end{itemize}

\section{Input file}

The parameter file \texttt{ics.nml} for a Larson fractal:
\begin{lstlisting}
&ics
  outfile = 'data/frac1k.dat'
  n       = 1000
  seed    = 12345
  distribution = 'fractal'
  radius  = 1.0            ! total radius (pc)
/
&fractal
  fdim = 1.6
  velocities = 'larson'
  vlarson = 1.0            ! km/s at 1 pc
/
\end{lstlisting}
Run it in \texttt{HENDOS/run} with \texttt{./makeics ics.nml}. All the
inputs are described in the file \texttt{PARAMS\_README.md} in that
directory.

\section*{References}
\begin{itemize}\setlength\itemsep{0pt}
\item Goodwin S.~P., Whitworth A.~P., 2004, A\&A, 413, 929
\item Heyer M., Krawczyk C., Duval J., Jackson J.~M., 2009, ApJ, 699,
  1092
\item Larson R.~B., 1981, MNRAS, 194, 809
\item Schoenberg I.~J., 1938, Trans.\ Amer.\ Math.\ Soc., 44, 522
\item Solomon P.~M., Rivolo A.~R., Barrett J., Yahil A., 1987, ApJ, 319,
  730
\end{itemize}

\end{document}
